\section{Extended Analysis: Strategic Pressure, Theory, and Scope}

\paragraph{Strategic pressure on evaluators.}
When benchmarks become optimization targets, evaluators become attack surfaces.
In coding and formal-reasoning domains, agents can discover rubric-sensitive shortcuts that improve judged scores without improving semantic correctness.
This failure mode is familiar in LLM-as-a-judge studies and in broader benchmark-gaming discussions.
Our certification protocol treats this risk explicitly by requiring directional behavioral checks under controlled degradations.
If a judge cannot reliably rank easier versus harder perturbation states, then it should not be trusted as the sole optimization objective for agent training loops.

\paragraph{Why monotonicity checks matter in practice.}
Monotonicity is a low-level but high-leverage validity criterion.
If an evaluator scores a more heavily bugged artifact above a less bugged one, optimization dynamics can amplify the error.
In iterative agent loops, this can induce selection for stylistic artifacts that correlate with judge preferences while breaking intended behavior.
Bug and specification ladders are therefore not only diagnostics.
They are preventive tests against reward misspecification in evaluator-mediated training and search.

\paragraph{Connection to interactive measurement.}
Interactive evaluation differs from static benchmark scoring because an evaluator's signal changes the policy trajectory.
A weak evaluator can create feedback loops where successive proposals become increasingly tuned to evaluator idiosyncrasies.
A certified evaluator does not eliminate this risk, but it raises the bar by enforcing auditable invariants at the score-function level.
This makes iterative improvement pipelines easier to monitor and easier to shut down when drift indicators fire.

\paragraph{Formal-methods framing and evaluator role.}
In Lean-oriented workflows, kernel checking verifies whether a submitted proof is valid under a stated theorem.
The evaluator's role is orthogonal: assessing whether the stated theorem/specification is aligned with task intent.
This distinction matters for measurement claims.
Kernel-level soundness does not imply specification quality, and specification quality cannot be reduced to syntax-level matching.
Judge certification therefore targets the semantic layer between raw generated artifact and intended task meaning.

\paragraph{Specification-relative interpretation of scores.}
Judge scores are specification-relative by design.
A high score means strong conformance to the chosen reference and rubric, not universal correctness across all possible properties.
This is not a weakness unique to LLM judges.
It is a general limitation of finite evaluation under open-ended semantics.
Making this conditionality explicit improves scientific communication and reduces over-interpretation of benchmark numbers.

\paragraph{Residual validity risks.}
Even when P1--P4 are strong, several risks remain:
\begin{itemize}
\item \textbf{content validity gaps}: perturbation ladders may not cover all error classes observed in deployment;
\item \textbf{criterion drift}: correlation with humans can degrade as rubric standards evolve;
\item \textbf{construct shift}: new domains may require new notions of theorem quality;
\item \textbf{optimization leakage}: agents may overfit to known perturbation families.
\end{itemize}
For this reason, we view trust indices as dynamic artifacts with explicit maintenance requirements.

\paragraph{Operational recommendations.}
Based on our experiments, we recommend four implementation rules for teams deploying judge-based evaluation:
\begin{enumerate}
\item publish per-property scores alongside aggregate trust values;
\item maintain a frozen anchor subset for longitudinal reassessment;
\item run disagreement audits between primary and secondary judges;
\item treat large prompt or model updates as mandatory re-certification events.
\end{enumerate}
These rules keep evaluator governance lightweight while preserving interpretability under change.

\paragraph{Relation to robustness-without-ground-truth methods.}
Information-theoretic evaluation approaches and incentive-compatible query designs offer complementary tools in settings with weak or absent gold references.
Our setting includes curated references, so we emphasize supervised property checks.
A promising direction is hybridization: combine robust unsupervised disagreement signals with supervised perturbation diagnostics and human-anchored assessment.
Such hybrids may improve resilience to evaluator gaming while retaining semantic grounding.

\paragraph{Why this scope is suitable for a workshop contribution.}
This paper does not claim a final universal metric for formal-verification judges.
Instead, it contributes a practical, testable methodology and an empirical case study showing that the methodology can support useful ranking decisions.
For AI Measurement Science, this is the key step: move from evaluator intuition to evaluator evidence.
The broader research program is to standardize these evidence bundles across tasks, domains, and model generations.

\paragraph{A minimal reporting template for evaluator cards.}
To support comparability across papers, we propose a concise evaluator-card template with five fields.
Field 1 is \emph{intended use}: rank model variants for research iteration, not certify deployment correctness.
Field 2 is \emph{evidence}: development panel size, rating protocol, and measured human-alignment statistics.
Field 3 is \emph{behavioral validity}: per-property scores with confidence intervals where possible.
Field 4 is \emph{known blind spots}: uncovered perturbation classes and domains with weak support.
Field 5 is \emph{maintenance policy}: re-certification triggers and latest assessment date.
This template is short enough for workshop papers and strong enough to improve reproducibility.

\paragraph{Error analysis priorities for next iterations.}
Our current protocol highlights whether monotonicity holds, but not always \emph{why} failures occur.
The next iteration should classify monotonicity violations by semantic category:
arithmetic specification drift, boundary-condition omission, theorem weakening, and proof-irrelevant style changes.
That taxonomy would make judge-assessment interventions more targeted.
For example, if failures cluster in boundary-condition omissions, perturbation suites should over-sample those edits and judge prompts should include boundary-sensitive criteria.
This turns property failures into actionable engineering signals.

\paragraph{Human factors in measurement operations.}
Evaluator quality is partly a human-systems problem.
Annotator instructions, rubric granularity, and quality-control policies can all alter apparent judge alignment.
Even with fixed models, small rubric wording changes can shift correlations enough to reorder judge candidates.
For that reason, judge-assessment metadata should be versioned as carefully as model checkpoints.
We recommend including rubric hashes and annotation-guideline versions in experiment logs to keep judge comparisons auditable over time.

\paragraph{Positioning against ``single-number'' evaluation culture.}
A recurring anti-pattern in LLM evaluation is collapsing complex reliability behavior into one headline number.
Single-number reporting is attractive for leaderboard communication but often hides the mechanisms behind failure.
Property-based certification addresses this by preserving a decomposition:
identity, directional monotonicity, and stability are reported separately and then aggregated for selection.
The aggregate is still useful, but it is no longer opaque.
For workshop goals focused on rigorous evaluation, this decomposition is a practical compromise between scientific clarity and operational simplicity.

\paragraph{Takeaway for AI measurement science.}
As AI systems become increasingly agentic, evaluator design must be treated as model design.
A judge is not only a metric; it is part of the control loop that shapes what systems learn to optimize.
CertJudge provides one concrete recipe for this perspective:
define falsifiable evaluator properties, measure agreement with humans once, monitor under drift, and re-certify on change.
We view this as a broadly reusable pattern for moving from benchmark-centric evaluation to measurement-centric evaluation.
